%\documentclass{book}\usepackage{sjtfonts}\usepackage{includegraphics}\usepackage{alltt}\usepackage{latexsym}
%\usepackage{array}\usepackage{multicol}
%\begin{document}\input{defs0}%		miradefs.tex 		    June 1994

% opening and closing a quoted Miranda section

\newcommand{\mira}{\noindent\hrulefill\nopagebreak\so}
\newcommand{\arim}{\nopagebreak\st\hrulefill\\}

% backslash in the Miranda environment
% Only feasible approach is to use \verb+\+
% in line!

%\newcommand{\bs}{$\mi\backslash\!\!$}
%\newcommand{\lbs}{\mbox{\verb+\+}}

% ignoring part of the input

\newcommand{\ignore}[1]{}

% a blank line in a Miranda section

\newcommand{\bl}{\mbox{\ }}

% Signalling a diffculty.

\definecolor{light-gray}{gray}{0.95}

\newcommand{\beware}[2]
 {\medskip\noindent\colorbox{light-gray}{\parbox{\textwidth}
 {\mbox{\large\textbf{#1}} \medskip\noindent\\ #2}\medskip\noindent}}

% Subscripting in Miranda text
\newcommand{\subscr}[2]{\mbox{${\mbox{\texttt{#1}}}_{\mbox{\texttt{#2}}}$}}

%\newcommand{\subscr}[2]{\mbox{${\mbox{\mi #1}}_{\mbox{\smtt #2}}$}}
\newcommand{\aone}{\subscr{a}{1}}
\newcommand{\atwo}{\subscr{a}{2}}
\newcommand{\ak}{\subscr{a}{k}}

\newcommand{\aoneone}{\subscr{a}{1,1}}

\newcommand{\eone}{\subscr{e}{1}}
\newcommand{\etwo}{\subscr{e}{2}}
\newcommand{\ei}{\subscr{e}{i}}
\newcommand{\er}{\subscr{e}{r}}

\newcommand{\fone}{\subscr{f}{1}}
\newcommand{\fl}{\subscr{f}{l}}

\newcommand{\gone}{\subscr{g}{1}}
\newcommand{\gtwo}{\subscr{g}{2}}
\newcommand{\gi}{\subscr{g}{i}}
\newcommand{\gr}{\subscr{g}{r}}

\newcommand{\hone}{\subscr{h}{1}}
\newcommand{\hl}{\subscr{h}{l}}

\newcommand{\pone}{\subscr{p}{1}}
\newcommand{\ptwo}{\subscr{p}{2}}
\newcommand{\pk}{\subscr{p}{k}}

\newcommand{\qone}{\subscr{q}{1}}
\newcommand{\qtwo}{\subscr{q}{2}}
\newcommand{\qk}{\subscr{q}{k}}

\newcommand{\rone}{\subscr{r}{1}}
\newcommand{\rtwo}{\subscr{r}{2}}

\newcommand{\tone}{\subscr{t}{1}}
\newcommand{\ttwo}{\subscr{t}{2}}
\newcommand{\tn}{\subscr{t}{n}}

\newcommand{\vone}{\subscr{v}{1}}
\newcommand{\vtwo}{\subscr{v}{2}}
\newcommand{\vn}{\subscr{v}{n}}

% for regular expressions

\newcommand{\stone}{\subscr{st}{1}}
\newcommand{\sttwo}{\subscr{st}{2}}
\newcommand{\stn}{\subscr{st}{n}}
\newcommand{\rthree}{\subscr{r}{3}}

\newcommand{\eps}{$\varepsilon$}
\newcommand{\trans}[1]{\stackrel{\mbox{\tt #1}}{\longrightarrow}}



% superscript in Miranda text

\newcommand{\superscr}[2]{\mbox{${\mbox{\mi #1}}^{\mbox{\smtt #2}}$}}

% Not equals in Miranda text

\newcommand{\noteq}{\makebox[0pt][l]{=}/}
\newcommand{\overslash}{\makebox[0pt][l]{/}}

% Definitions in complexity chapter

\newfont{\oldSym}{cmsy10}
\newcommand{\Th}{\boldmath$\oldSym\Theta$}
\newcommand{\pow}[1]{\^{}#1}
\newcommand{\leqq}{$\leq$}
\newcommand{\geqq}{$\geq$}
\newcommand{\lll}{$\ll$}
\newcommand{\eqq}{$\equiv$}
\newcommand{\ltwo}{\subscr{log}{2}}

% Indexing

\newcommand{\minx}[1]{\index{#1@{\ttfamily{}#1}}}
\appendix


\newcommand{\gitem}[1]{\medskip\noindent{\bf #1\ \ \ }
                       %\raggedright
                       }

\newcommand{\upp}[1]{\ensuremath{{}^{\mbox{\footnotesize\texttt #1}}}}

%\minitocoptionfalse
\chapter{Glossary}

We include this glossary to give a quick reference to the most
widely used terminology in the book. Words appearing in \textbf{bold} in the
descriptions have their own entries. Further references and examples are
to be found by consulting the index.
\normalfont\normalsize
\hspace*{12pt} 


\begin{multicols}{2}

{\raggedright

\gitem{Abstract type}\index{abstract data type}
An abstract type definition consists of the type name, the \textbf{signature} of
the type, and the implementation equations for the names in the
signature.

\gitem{Algebraic type}\index{algebraic type}
An algebraic type definition states what are the \textbf{constructors} of the
type. For instance, the declaration
\begin{alltt}
data Tree = Leaf Int | 
            Node Tree Tree
\end{alltt}
says that the two constructors of the \texttt{Tree} type are
\texttt{Leaf} and \texttt{Node}, and that their types are, respectively,
\begin{alltt}
Leaf :: Int->Tree
Node :: Tree->Tree->Tree
\end{alltt}

\gitem{Application}\index{function application}
This means giving values to (some of) the arguments
of a function. If an \texttt{n}-argument function is given fewer than \texttt{n}
arguments, this is called a \textbf{partial application}.
Application is written using \textbf{juxtaposition}.

\gitem{Argument}\index{argument}
A \textbf{function} takes one or more arguments into an \textbf{output}. Arguments
are also known as \textbf{inputs} and \textbf{parameters}.

\gitem{Associativity}\index{associativity}
The way in which an expression involving two applications of an operator is
interpreted. If \texttt{x\#y\#z} is interpreted as \texttt{(x\#y)\#z} then
\texttt{\#} is left associative, if as \texttt{x\#(y\#z)} it is right associative;
if both bracketings give the same result then \texttt{\#} is called associative.

\gitem{Base types}\index{type!base}
The types of numbers, including
\texttt{Int} and \texttt{Float}, \textbf{Booleans}, \texttt{Bool}, and \textbf{characters},
\texttt{Char}.

\gitem{Binding power}\index{operator!binding power of}
The `stickiness' of an operator, expressed as an integer; the higher the
number the stickier the operator. For example, \texttt{2+3*4} is interpreted as
\texttt{2+(3*4)} as `\texttt{*}' has higher binding power -- binds more tightly 
-- than `\texttt{+}'.

\gitem{Booleans}\index{Booleans}
The type containing the two `truth values' \texttt{True} and \texttt{False}.

\gitem{Calculation}\index{calculation}
A calculation is a line-by-line \textbf{evaluation} of a Haskell \textbf{expression}
on paper. Calculations use the \textbf{definitions} which are contained in a
\textbf{script} as well as the built-in definitions.

\gitem{Cancellation}\index{type checking!rule of cancellation}
The rule for finding the type of a partial application.

\gitem{Character}\index{characters}
A single letter, such as \texttt{'s'} or \texttt{'\bs{t}'}, the tab character.
They form the \texttt{Char} type.

\gitem{Class}\index{classes}
A collection of types. A class is defined by specifying a \textbf{signature};
a type is made an \textbf{instance} of the class by supplying an
implementation of the definitions of the signature over the type.

\gitem{Clause}\index{clause}
A clause is one of the alternatives making up a \textbf{conditional 
equation}. A clause consists of a \textbf{guard} followed by an \textbf{expression}. 
When evaluating a function application, the first clause whose
guard evaluates to \texttt{True} is chosen.

\gitem{Combinator}
Another name for a \textbf{function}.

\gitem{Comment}\index{comment}
Part of a \textbf{script} which plays no computational role; it is there for the
reader to read and observe. Comments are specified in two ways: the part of
the line to the right is made a comment by the symbol \texttt{--}; a comment of
arbitrary length is enclosed by \texttt{\{-} and \texttt{-\}}.

\gitem{Complexity}\index{complexity}
A measurement of the time or space behaviour of a function.

\gitem{Composition}\index{function composition}
The combination of two functions by passing the \textbf{output} of one to the
\textbf{input} of the other.

\gitem{Concatenate}\index{concatenate}
To put together a number of lists into a single list.

\gitem{Conditional equation}\index{conditional equation}
A conditional equation consists of a left-hand side followed by a number of
\textbf{clauses}. Each clause consists of a \textbf{guard} followed by an expression
which is to be equated with the left-hand side of the \textbf{equation} if that
particular clause is chosen during evaluation. The clause chosen is the first
whose guard evaluates to \texttt{True}.


\gitem{Conformal pattern match}\index{pattern matching!conformal}
An equation in which a pattern appears on the left-hand side of an equation,
as in
\begin{alltt}
(x,y) = ....
\end{alltt}

\gitem{Constructor}\index{constructor}
An \textbf{algebraic type} is specified by its constructors, which are the
functions which build elements of the algebraic type. 
\itempar
In the example in the
entry for algebraic types, elements of the type are constructed using 
\texttt{Leaf} and \texttt{Node}; the elements are \texttt{Leaf n} where \texttt{n::Int} and
\texttt{Node s t} where \texttt{s} and \texttt{t} are trees.

\gitem{Context}\index{context}
The hypotheses which appear before \texttt{=>} in type and class declarations. A
context \texttt{M a} means that the type \texttt{a} must belong to the class \texttt{M}
for the function or class definition to apply. For instance, to apply a
function of type
\begin{alltt}
Eq a => [a] -> a -> Bool
\end{alltt}
to a list and object, these must come from types over which equality is
defined.

\gitem{Curried function} \index{currying}
A function of at least two arguments which takes its arguments one at a time, so
having
the type
\begin{alltt}
t1 -> t2 -> ... -> t
\end{alltt}
in contrast to the {\em uncurried\/} version
\begin{alltt}
(t1,t2,...) -> t
\end{alltt}
The name is in honour of Haskell B.\ Curry, after whom the Haskell language is
also named.

\gitem{Declaration}\index{declaration}
A \textbf{definition} can be accompanied by a statement of the \textbf{type} of
the object defined; these are often called type declarations.  

\gitem{Default}\index{default}
A default holds in the absence of any other definition. Used in \texttt{class}
definitions to give  definitions of some of the operations in
terms of others; an example is the definition of \texttt{/=} in
the \texttt{Eq} class.

\gitem{Definition}\index{definition}
A definition associates a \textbf{value} or a \textbf{type} with a \textbf{name}.

\gitem{Design}\index{design}
In writing a system, the effort expended {\em before\/} implementation is
started.

\gitem{Derived class instance}\index{instance!derived}
An instance of a standard class which is derived by the system, rather than
put in explicitly by the programmer. 

\gitem{Enumerated type}\index{enumerated type}
An \textbf{algebraic type} with each constructor having no arguments.

\gitem{Equation}\index{equation}\index{equation!conditional}
A \textbf{definition} in Haskell consists of a number of equations. On the left-hand side of the equation is a \textbf{name} applied to zero or more 
\textbf{patterns}; on the right-hand side is a value. In many
cases the equation is \textbf{conditional} and has two or more
\textbf{clauses}. Where the meaning is clear we shall sometimes
take `equation' as shorthand for `equation or conditional
equation'.

\gitem{Evaluation}\index{evaluation}
Every \textbf{expression} in Haskell has a value; evaluation is the process of
finding that value. A \textbf{calculation} evaluates an expression, as does
an interactive
Haskell system when that expression is typed to the prompt.

\gitem{Export}\index{export}
The process of defining which definitions will be visible when a module is
\textbf{imported} by another.

\gitem{Expression}\index{expression}
An expression is formed by applying a \textbf{function} or \textbf{operator} to its
arguments; these arguments can be \textbf{literal} values, or expressions
themselves. A simple numerical expression is 
\begin{alltt}
(2+8)-10
\end{alltt}
in which the
operator `\texttt{-}' is applied to two arguments.

\gitem{Extensionality}\index{extensionality principle}
The principle of proof which says that two functions are equal if they give
equal results for every input.

\gitem{Filter}\index{filter}
To pick out those elements of a list which have a particular property,
represented by a \textbf{Boolean}-valued function.

\gitem{Floating-point number}\index{numbers!floating point}
A number which is given in decimal (e.g.\ {\mi456.23}) or exponent (e.g.\
\texttt{4.5623e+2}) form; these numbers form the type \texttt{Float}.

\gitem{Fold}\index{folding}
To combine the elements of a list using a binary \textbf{operation}.

\gitem{Forward composition}\index{forward composition}
Used for the operator `\texttt{>.>}' with the definition
\begin{alltt}
f >.> g = g . f
\end{alltt}
\texttt{f >.> g} can be read `\texttt{f} then \texttt{g}'.

\gitem{Function}\index{function}
A function is an object which returns a \textbf{value}, called the 
\textbf{output} or \textbf{result} when it is applied to its \textbf{inputs}. The
inputs are also known as its \textbf{parameters} or \textbf{arguments}. 
 
\noindent
Examples
include the square root function, whose input and output are numbers, and the
function which returns the borrowers (output) of a book (input) in a database
(input).

\gitem{Function types}
The type of a \textbf{function} is a function type, so that, for instance, the
function which checks whether its integer
argument is even has type \texttt{Int->Bool}.
This is the type of functions with \textbf{input} type \texttt{Int} and 
\textbf{output} type \texttt{Bool}.

\gitem{Generalization}\index{generalization}
Replacing an object by something of which the original object is an instance.
\itempar
This might be the replacement of a function by a polymorphic function
from which the original is obtained by passing the appropriate parameter, or
replacing a logical formula by one which implies the original.

\gitem{Guard}\index{guard}
The \textbf{Boolean} expression appearing to the right of `\texttt{|}' and to the left
of `\texttt{=}' in a 
\textbf{clause} of a \textbf{conditional equation} in a Haskell
\textbf{definition}.

\gitem{Higher-order function}\index{higher-order function}
A \textbf{function} is higher-order if either one of its \textbf{arguments} or
its \textbf{result}, or both, are functions.

\gitem{Identifier} \index{identifier}
Another word for \textbf{name}.

\gitem{Implementation}
The particular \textbf{definitions} which make a design concrete; for an 
\textbf{abstract data type}, the definitions of the objects named in the 
\textbf{signature}.

\gitem{Import}\minx{import}
The process of including the \textbf{exported} definitions of one module in
another.

\gitem{Induction}\index{induction}
The name for a collection of methods of proof, by which statements of the
form `for all \texttt{x} \ldots' are proved.

\gitem{Infix}\index{function!infix}
An \textbf{operation} which appears between its \textbf{arguments}. Infix functions
are called \textbf{operators}.

\gitem{Inheritance}\index{inheritance}
One \textbf{class} inherits the operations of another if the first class is in
the \textbf{context} of the definition of the second. For instance, of the
standard classes, \texttt{Ord} inherits (in)equality from \texttt{Eq}.

\gitem{Input}\index{input}
A \textbf{function} takes one or more inputs into an \textbf{output}. Inputs
are also known as \textbf{arguments} and \textbf{parameters}. The `square'
function takes a single numerical input, for instance.

\gitem{Instance}\index{instance}\index{instance!of variable}
\index{instance!of class}
The term `instance' is used in two different ways in Haskell.
\itempar
An instance of a \textbf{type}
is a type which is given by \textbf{substituting} a type expression for a
type \textbf{variable}. For example, \texttt{[(Bool,b)]} is an instance of 
\texttt{[a]}, given by substituting the type
\texttt{(Bool,b)} for the variable \texttt{a}.
\itempar
An instance of a \textbf{class}, such as \texttt{Eq (a,b)},
is given by declaring how the function(s) of the class, in this case \texttt{==},
are defined over the given type (here \texttt{(a,b)}). Here we would say
\begin{alltt}
(x,y) == (z,w)
  = (x==z) \&\& (y==w)
\end{alltt}

\gitem{Integers}\index{numbers!integers}
The positive and negative whole numbers. In Haskell the type \texttt{Int}
represents the integers in a fixed size, while the type \texttt{Integer}
represents them 
exactly, so that evaluating \texttt{2} to the power \texttt{1000} will give a
result consisting of some three hundred digits.

\gitem{Interactive program}
A program which reads from and writes to the terminal; reading and writing
will be {\em interleaved}, in general.

\gitem{Interface}\index{interface}
The common information which is shared between two program modules.

\gitem{Juxtaposition}\index{juxtaposition}
Putting one thing next to another; this is the way in which function
application is written down in Haskell.

\gitem{Lambda expression}\index{lambda expression}
An \textbf{expression} which denotes a \textbf{function}. After a `\texttt{\bs{}}' we
list the arguments of the function, then an `\texttt{->}' and then the result.
For instance, to add a number to the length of a list we could write
\begin{alltt}
\bs{xs} n -> length xs + n
\end{alltt}
The term `lambda' is used since `\texttt{\bs{}}' is close to the Greek letter
`$\lambda$', or lambda, which is used in a similar way in Church's
lambda calculus.

\gitem{Lazy evaluation}\index{lazy evaluation}
The sort of expression \textbf{evaluation} in Haskell. In a function
application only those arguments whose values are {\em needed\/} will be
evaluated, and moreover, only the parts of structures which are needed
will be examined.

\gitem{Linear complexity}\index{complexity!linear}
Order \texttt{1}, \texttt{O(n\upp{1})}, behaviour.

\gitem{Lists}\index{lists}
A list consists of a collection of elements of a particular type, given in
some order, potentially containing a particular item more than once.
The list \texttt{[2,1,3,2]} is of type \texttt{[Int]}, for example.

\gitem{Literal}\index{literal}
Something that is `literally' a value: it needs no \textbf{evaluation}. Examples
include \texttt{34}, \texttt{[23]} and \texttt{"string"}.

\gitem{Local definitions}\index{local definitions}
The definitions appearing in a \texttt{where} clause or a \texttt{let} expression.
Their \textbf{scope} is
the equation or expression to which the clause or \texttt{let} is attached.

\gitem{Map}\index{mapping} To apply an operation to every element of a list.

\gitem{Mathematical induction}\index{mathematical induction}
A method of proof for statements of the form `for all natural numbers 
\texttt{n}, the statement \texttt{P(n)} holds'. 
\itempar
The proof is in two parts: the
base case, at zero, and the induction step, at which \texttt{P(n)} is proved
on the assumption that \texttt{P(n-1)} holds.

\gitem{Memoization}\index{memoization}
Keeping the value of a sub-computation (in a list, say)
so that it can be reused
rather than recomputed, when it is needed.

\gitem{Module}\index{modules}
Another name for a \textbf{script}; used particularly when more than one script
is used to build a program.

\gitem{Monad}\index{monad}
A monad consists of a type with (at least) two functions, \texttt{return} and
\texttt{>>=}. Informally, a monad can be seen as performing some sorts of action
before returning an object. The two monad functions
respectively return a value without any action, and sequence two
monadic operations.

\gitem{Monomorphic} \index{type!monomorphic}
A type is \textbf{monomorphic} if it is not \textbf{polymorphic}.

\gitem{Most general type}\index{type!most general}
The most general type of an expression is the type \texttt{t} with the property
that every other type for the expression is an \textbf{instance} of \texttt{t}.

\gitem{Mutual recursion}\index{recursion!mutual}
Two definitions, each of which depends upon the other.

\gitem{Name}\index{name}
A \textbf{definition} associates a name or \textbf{identifier} with a value. Names
of \textbf{classes}, \textbf{constructors} and \textbf{types}
must begin with capital letters; names of 
\textbf{values}, \textbf{variables} and \textbf{type variables}
begin with small letters. After the first letter, any
letter, digit, `\texttt{'}' or `\texttt{\_}' can be used.

\gitem{Natural numbers}\index{numbers!natural numbers}
The non-negative whole numbers: \texttt{0}, \texttt{1}, \texttt{2}, \ldots\ .

\gitem{Offside rule}\index{offside rule}
The way in which the end of a part of a definition is expressed using the
{\em layout\/} of a \textbf{script}, rather than an explicit symbol for the end.

\gitem{Operation}
Another name for \textbf{function}.

\gitem{Operator}\index{operator}
A \textbf{function} which is written in infix form, between its \textbf{arguments}.
The function \texttt{f} is made infix thus: \texttt{`f`}. 

\gitem{Operator section}\index{operator sections}
A partially applied operator.

\gitem{Output}\index{output}
When a \textbf{function} is applied to one or more \textbf{inputs}, the resulting
value is called the output, or \textbf{result}. Applying the `square' function
to \texttt{(-2)} gives the output \texttt{4}, for example.

\gitem{Overloading}\index{overloading}
The use of the same \textbf{name} to mean two (or more) different things, at
different types. The equality operation, \texttt{==},  is an example. Overloading
is
supported in Haskell by the \textbf{class} mechanism.

\gitem{Parameter}\index{parameter}
A \textbf{function} takes one or more parameters into an \textbf{output}. Parameters
are also known as \textbf{arguments} and \textbf{inputs}, and applying a function
to its inputs is sometimes known as `passing its parameters'. 

\gitem{Parsing}\index{parsing}
Revealing the structure of a sentence in a formal language.

\gitem{Partial application}\index{partial application}
A \textbf{function} of type \texttt{\tone->\ttwo->\ldots->\tn->t} can
be applied to \texttt{n} arguments, or less. In the latter case, the
\textbf{application} is partial, since the result can itself be passed further
parameters.

\gitem{Pattern}\index{pattern}
A pattern is either a \textbf{variable}, a \textbf{literal}, a \textbf{wild card}
or the application of a \textbf{constructor}
to other patterns. 
\itempar
The term `pattern' is also used as short for a `pattern of
computation' such as `applying an operation to every member of
a list', a pattern which in Haskell is realised by the
\texttt{map} function.

\gitem{Polymorphism} \index{polymorphism}
A type is polymorphic if it contains type \textbf{variables}; such a type will
have many \textbf{instances}.

\gitem{Prefix}\index{function!prefix}
An \textbf{operation} which appears before its \textbf{arguments}.

\gitem{Primitive recursion}\index{primitive recursion}
Over the natural numbers, defining the values of a function outright at zero,
and at \texttt{n} greater than zero using the value at \texttt{n-1}. 

\noindent
Over an \textbf{algebraic type}
defining the function by cases over the constructors; recursion is permitted
at arguments to a constructor which are of the type in question.

\gitem{Proof}\index{proof}
A logical argument which leads us to accept a logical statement as being
valid.

\gitem{Pure programming language}\index{functional programming!pure}
A functional programming
language is pure if it does not allow \textbf{side-effects}.

\gitem{Quadratic complexity}\index{complexity!quadratic}
Order two, \texttt{O(n\upp{2})}, behaviour.

\gitem{Recursion}\index{recursion}
Using the name of a value or type in its own \textbf{definition}.

\gitem{Result}\index{result}
When a \textbf{function} is applied to one or more \textbf{inputs}, the resulting
value is called the result, or \textbf{output}.

\gitem{Scope}\index{scope}
The area of a program in which a \textbf{definition} or definitions are
applicable. 

\noindent
In Haskell the scope of top-level definitions is by default the
whole \textbf{script} in which they appear; it may be extended by importing the
module into another.
More limited scopes are given by \textbf{local
definitions}.

\gitem{Script}\index{script}
A script is a file containing \textbf{definitions}, \textbf{declarations} and
module statements.

\gitem{Set}\index{sets}
A collection of objects for which the order of elements and the number of
occurrences of each element are irrelevant.

\gitem{Side-effect}\index{side-effect}
In a language like Pascal, evaluating an expression can cause other things to
happen besides a value being computed. These might be I/O operations, or
changes in values stored. In Haskell this does not happen, but a \textbf{monad}
can be used to give a similar effect, without compromising the simple model of
evaluation underlying the language. Examples are \texttt{IO} and \texttt{State}.

\gitem{Signature}\index{signature}
A sequence of type \textbf{declarations}. These declarations state what are the
types of the operations (or functions) over an \textbf{abstract type} or a 
\textbf{class} which can be
used to manipulate elements of that type.

\gitem{Stream}\index{stream}
A stream is a channel upon which items arrive in sequence; in Haskell we can
think of \textbf{lazy} lists in this way, so it becomes a synonym for lazy list.

\gitem{String}\index{strings}
The type \texttt{String} is a \textbf{synonym} for lists of characters, \texttt{[Char]}.

\gitem{Structural induction}\index{structural induction}
A method of proof for statements of the form `for all
finite lists 
\texttt{xs}, the statement \texttt{P(xs)} holds of \texttt{xs}'. The proof is in two parts: the
base case, at \texttt{[]}, and the induction step, at which \texttt{P(y:ys)} is proved
on the assumption that \texttt{P(ys)} holds. 

\noindent
Also used of the related principle
for any algebraic type.

\gitem{Substitution}\index{substitution}
The replacement of a \textbf{variable} by an \textbf{expression}. For example, 
\texttt{(9+12)} is given by substituting \texttt{12} for \texttt{n} in \texttt{(9+n)}.
Types can also be substituted for type variables; see the entry for 
\textbf{instance}.

\gitem{Synonym}\index{type synonym}
Naming a type is called a type synonym. The keyword \texttt{type} is used for
synonyms.

\gitem{Syntax}\index{syntax}
The description of the properly formed programs (or sentences) of a language.

\gitem{Transformation}\index{program transformation}
Turning one program into another program which computes identical results,
but with different behaviour in other respects such as time or space
efficiency.

\gitem{Tuples}\index{tuples}
A tuple type is built up from a number of component types. Elements of the
type consist of tuples of elements of the component types, so that
\begin{alltt}
(2,True,3) :: (Int,Bool,Int)
\end{alltt}
for instance.

\gitem{Type}\index{type}
A collection of values. Types can be built from the \textbf{base} types
using \textbf{tuple}, \textbf{list} and \textbf{function types}. New types can be
defined using the \textbf{algebraic} and \textbf{abstract} type mechanisms, and
types can be named using the type \textbf{synonym} mechanism.

\gitem{Type variable}\index{type variable}
A \textbf{variable}
which appears in a \textbf{polymorphic type}. An \textbf{identifier} beginning with a
small letter can be used as a type variable; in this text we use the
letters at the
start of the alphabet,
\texttt{a}, \texttt{b}, \texttt{c} and so on.

\gitem{Undefinedness}\index{undefinedness}
The result of an expression whose evaluation continues forever, rather than
giving a {\em defined\/} result.

\gitem{Unification}\index{unification}
The process of finding a common \textbf{instance} of two (type) expressions
containing (type) variables.

\gitem{Value}\index{value}
A value is a member of some \textbf{type}; the value of an \textbf{expression} is the
result of \textbf{evaluating} the expression.

\gitem{Variable}\index{variable}
A variable stands for an {\em arbitrary\/} value, or in the case of type
variables, an arbitrary type. Variables and type variables
have the same syntax as \textbf{names}.

\gitem{Verification}
Proving that a function or functions have particular logical properties.

\gitem{Where clause}\minx{where}
Definitions \textbf{local} to a (conditional) \textbf{equation}.

\gitem{Wild card}\index{pattern!wild card}
The name for the pattern `\texttt{\_}', which is matched by any value of the
appropriate type.

} % end raggedright

\end{multicols}


%\end{document}
